
\subsection{统一反冲参数与模型选择}\label{subsec:recoil-regime-selection-dp47}

单频近场交换与周期势衍射都把电子连接到相隔固定纵向动量步长的通道，因此是否可忽略反冲、是否可取均匀能量梯以及是否能约化到少数通道，都由同一组无量纲尺度比决定。对中心动量 $p_0$、纵向群速度 $v_0$ 和纵向有效质量
\begin{equation*}
 m_\parallel=\frac{1}{\partial_p^2E|_{p_0}},
\end{equation*}
取单步纵向动量 $\hbar\kappa_z$，定义统一反冲频率
\begin{equation}
\boxed{\omega_{\rm r}=\frac{\hbar\kappa_z^2}{2m_\parallel}.}
\label{eq:unified_recoil_frequency_dp47}
\end{equation}
二阶色散首先决定相邻反冲边的失谐。由 $\mathcal E_\ell=E(p_0+\ell\hbar\kappa_z)$ 的二阶展开，
\begin{align*}
\mathcal E_\ell
&=E_0+\ell\hbar v_0\kappa_z+\ell^2\hbar\omega_{\rm r},\\
\frac{(\mathcal E_{\ell+1}-\mathcal E_\ell)-\hbar v_0\kappa_z}{\hbar}
&=[(\ell+1)^2-\ell^2]\omega_{\rm r}
=(2\ell+1)\omega_{\rm r}.
\end{align*}
因此
\begin{equation}
\boxed{
\delta_\ell
\equiv\frac{(\mathcal E_{\ell+1}-\mathcal E_\ell)-\hbar v_0\kappa_z}{\hbar}
=(2\ell+1)\omega_{\rm r},
\qquad
|\delta_\ell|\le(2N_{\rm occ}+1)\omega_{\rm r}.}
\label{eq:recoil_adjacent_detuning_bound_dp68}
\end{equation}
对于占据 $|\ell|\le N_{\rm occ}$ 的反冲格，把这一最大失谐依次与作用时间、耦合带宽、可用动量窗和输入包络比较，得到四个控制量
\begin{equation}
\boxed{
\epsilon_T=(2N_{\rm occ}+1)\omega_{\rm r}T,
\quad
\epsilon_\Omega=\frac{(2N_{\rm occ}+1)\omega_{\rm r}}{\Omega_{\rm char}},
\quad
\epsilon_{\rm win}=\frac{N_{\rm occ}\hbar|\kappa_z|}{p_0},
\quad
\epsilon_{\rm env}=\frac{\sigma_p}{\hbar|\kappa_z|}.}
\label{eq:recoil_regime_vector_dp47}
\end{equation}
$\Omega_{\rm char}$ 是相关相邻格点的特征耦合速率，$\sigma_p$ 是电子输入纵向动量宽度。$\epsilon_T$ 判定有限时间内是否能分辨曲率；$\epsilon_\Omega$ 只回答“整个占据窗的反冲失谐是否与耦合带宽可比”；$\epsilon_{\rm win}$ 衡量占据动量窗相对中心动量的宽度；$\epsilon_{\rm env}$ 则区分单一反冲格和多个准动量纤维的平均。

若要进一步约化成 Bragg 型二态模型，还必须区分目标交换边与其余泄漏边。设外部驱动相对线性同步频率额外具有可调失谐 $\Delta_{\rm d}$，逐边裸失谐为
\begin{equation*}
\Delta_\ell^{(\rm d)}\equiv\delta_\ell-\Delta_{\rm d}.
\end{equation*}
把准备显式保留的目标通道投影记为 $P$，其余通道投影记为 $Q=1-P$，并定义被消去子空间的 Hamiltonian $H_Q\equiv QHQ$。真正控制消元的不是某一个裸对角元，而是目标能量 $E$ 到整个 $Q$ 子空间谱的距离
\begin{equation*}
\Delta_Q(E)
\equiv\frac1\hbar\operatorname{dist}\!\bigl(E,\operatorname{spec}H_Q\bigr).
\end{equation*}
因此一般的 Bragg 泄漏控制量定义为
\begin{equation}
\boxed{
\epsilon_{\rm Br,leak}(E)
\equiv\frac{\Omega_{\rm char}}{\Delta_Q(E)}.}
\label{eq:bragg_target_leakage_criterion_dp72}
\end{equation}
若 $Q$ 子空间内部耦合可忽略，$H_Q$ 在反冲通道基中近似对角，此时才退化为常用的
\begin{equation*}
\Delta_Q(E)\longrightarrow
\min_{\ell\notin\mathcal S_{\rm tar}}|\Delta_\ell^{(\rm d)}|.
\end{equation*}
二态选择还要求目标边位于所选驱动带宽内，例如把目标功率展宽窗明确定义为 $|\Delta_{\ell_*}^{(\rm d)}|\le\Omega_{\rm char}$；与此同时要求 $\epsilon_{\rm Br,leak}\ll1$。前者保证目标边能被驱动，后者保证整个被删子空间只产生小的占据和虚跃迁修正。这样“目标共振”和“非目标抑制”由两个不同的数学条件承担。

据此，经典 PINEM 的 Bessel 闭式需要 $\epsilon_T\ll1$；有限反冲且中心动量附近的单格数值模型适用于 $\epsilon_T\gtrsim1$、$\epsilon_{\rm win}\ll1$ 且三阶相位余项受控的区域；Bragg 二态模型还额外要求\eqnrefs{eq:bragg_target_leakage_criterion_dp72} 的目标边近共振和泄漏边远失谐同时成立。若 $\epsilon_{\rm env}\gtrsim1$，则必须对 Bloch 准动量或输入动量分布作平均，不能用单一中心动量代表整个电子束。

二阶格模型还需要检查被舍弃的首个色散项。把真实能量在 $p_0$ 附近展开到三阶，
\begin{equation}
\boxed{
 E(p_0+\ell\hbar\kappa_z)
 =E_0+\ell\hbar v_0\kappa_z+\ell^2\hbar\omega_{\rm r}
 +\frac{E_0^{(3)}}{6}(\ell\hbar\kappa_z)^3+O[(\ell\hbar\kappa_z)^4].}
\label{eq:recoil_cubic_taylor_dp68}
\end{equation}
在传播时间 $t$ 后，最高占据阶的三阶相位为
\begin{equation}
\boxed{
 \varphi_3^{\max}(t)
 =\frac{|E_0^{(3)}|}{6\hbar}
 (N_{\rm occ}\hbar|\kappa_z|)^3t .}
\label{eq:cubic_phase_error_dp47}
\end{equation}
当 $\varphi_3^{\max}\ll1$ 时二阶反冲格足以描述相位；这比只要求 $N_{\rm occ}\hbar|\kappa_z|\ll p_0$ 更直接，因为它同时包含实际传播时间。

\begin{derivation}[从\eqnrefs{eq:recoil_adjacent_detuning_bound_dp68}到\eqnrefs{eq:bragg_target_leakage_criterion_dp72}]
Bragg 二态约化的核心并不是“失谐看起来很大”，而是被删去通道的波函数分量能够被一个小范数严格控制。把目标子空间投影记为 $P$，其余反冲通道投影记为 $Q=1-P$。在驱动旋转框架中，把 Hamiltonian 分块写成
\begin{equation*}
H=
\begin{pmatrix}
H_P & V\\
V^\dagger & H_Q
\end{pmatrix},
\qquad
|\Psi\rangle=
\begin{pmatrix}|\psi_P\rangle\\|\psi_Q\rangle\end{pmatrix}.
\end{equation*}
对准能量（或定态本征值）$E$，方程 $(H-E)|\Psi\rangle=0$ 的 $Q$ 分量是
\begin{equation*}
(H_Q-E)|\psi_Q\rangle=-V^\dagger|\psi_P\rangle.
\end{equation*}
只要 $E$ 不落在被消去子空间的谱上，就可精确求解
\begin{equation*}
|\psi_Q\rangle
=-(H_Q-E)^{-1}V^\dagger|\psi_P\rangle.
\end{equation*}
因此由算符范数次乘性，
\begin{equation*}
\|\psi_Q\|
\le
\|(H_Q-E)^{-1}\|\,\|V\|\,\|\psi_P\|.
\end{equation*}
对自伴 $H_Q$，预解算符范数由谱定理精确给出
\begin{equation*}
\|(H_Q-E)^{-1}\|
=\frac1{\operatorname{dist}(E,\operatorname{spec}H_Q)}
=\frac1{\hbar\Delta_Q(E)}.
\end{equation*}
再以特征耦合速率定义 $\|V\|\le\hbar\Omega_{\rm char}$，便得到
\begin{equation*}
\frac{\|\psi_Q\|}{\|\psi_P\|}
\le
\frac{\Omega_{\rm char}}{\Delta_Q(E)}
\equiv\epsilon_{\rm Br,leak}(E).
\end{equation*}
这就是正文\eqnrefs{eq:bragg_target_leakage_criterion_dp72} 的一般形式。

常用的“最小逐边失谐”只是这个谱隙的对角近似。若
\begin{equation*}
H_Q=H_Q^{(0)}+W_Q,
\end{equation*}
其中 $H_Q^{(0)}$ 在反冲通道基中对角，且其相对目标能量的本征值为 $\hbar\Delta_\ell^{(\rm d)}$，则
\begin{equation*}
\Delta_{Q,0}(E)
=\min_{\ell\notin\mathcal S_{\rm tar}}|\Delta_\ell^{(\rm d)}|.
\end{equation*}
谱在有界扰动下的稳定性给出保守界
\begin{equation*}
\Delta_Q(E)
\ge
\Delta_{Q,0}(E)-\frac{\|W_Q\|}{\hbar}.
\end{equation*}
因此只有当 $\|W_Q\|/\hbar$ 相对裸谱隙可忽略，或已经显式从分母中扣除时，才可把正文的一般谱隙换成最小逐边裸失谐。

把消去解代回 $P$ 方程还得到精确 Schur 补
\begin{equation*}
\left[H_P-E-V(H_Q-E)^{-1}V^\dagger\right]|\psi_P\rangle=0,
\end{equation*}
所以非目标通道对目标动力学的首个虚跃迁修正具有量级界
\begin{equation*}
\|\delta H_P\|
\le
\frac{\hbar\Omega_{\rm char}^2}{\Delta_Q(E)}.
\end{equation*}
因此 $\epsilon_{\rm Br,leak}\ll1$ 同时控制被删通道的实际占据幅度和它们诱导的二阶能移。目标边本身则必须保留在 $P$ 空间，并单独满足所选功率展宽窗内的近共振条件；不能用同一个“大失谐”条件处理目标边和泄漏边。
\end{derivation}

\begin{derivation}[从\eqnrefs{eq:recoil_cubic_taylor_dp68}到\eqnrefs{eq:cubic_phase_error_dp47}]
写

\begin{align*}
E(p_0+\Delta p)
&=E_0+v_0\Delta p+\frac{(\Delta p)^2}{2m_\parallel}+\frac{E_0^{(3)}}{6}(\Delta p)^3+\cdots,\\
\Delta p&=\ell\hbar\kappa_z.
\end{align*}
三阶能量修正因此为
\begin{equation*}
\Delta E_\ell^{(3)}=\frac{E_0^{(3)}}{6}(\ell\hbar\kappa_z)^3.
\end{equation*}
对相对论色散 $E(p)=\sqrt{m_{\mathrm e}^2c^4+c^2p^2}$，逐次求导得到
\begin{align*}
E'(p)&=\frac{c^2p}{E},\\
E''(p)&=\frac{m_{\mathrm e}^2c^6}{E^3},\\
E'''(p)&=-\frac{3m_{\mathrm e}^2c^8p}{E^5}.
\end{align*}
在 $p_0=\gamma_0m_{\mathrm e} v_0$、$E_0=\gamma_0m_{\mathrm e} c^2$ 处，
\begin{equation*}
E_0^{(3)}=-\frac{3v_0}{\gamma_0^4m_{\mathrm e}^2c^2}.
\end{equation*}
传播时间 $t$ 把能量误差累积成相位误差
\begin{equation*}
\varphi_{3,\ell}(t)=\frac{\Delta E_\ell^{(3)}t}{\hbar}=\frac{E_0^{(3)}}{6\hbar}(\ell\hbar\kappa_z)^3t.
\end{equation*}
在 $|\ell|\le N_{\rm occ}$ 上取绝对值最大值，得到
\begin{equation*}
\varphi_3^{\max}(t)=\frac{|E_0^{(3)}|}{6\hbar}(N_{\rm occ}\hbar|\kappa_z|)^3t,
\end{equation*}
即\eqnrefs{eq:cubic_phase_error_dp47}。相应的二阶最大相位为
\begin{equation*}
\varphi_2^{\max}(t)
=\frac{(N_{\rm occ}\hbar|\kappa_z|)^2t}{2m_\parallel\hbar}.
\end{equation*}
两者之比把几何窗口参数与真实高阶色散联系起来：
\begin{align*}
\frac{\varphi_3^{\max}}{\varphi_2^{\max}}
&=\frac{|E_0^{(3)}|m_\parallel}{3}
N_{\rm occ}\hbar|\kappa_z|\\
&=\frac{v_0^2}{c^2}
\frac{N_{\rm occ}\hbar|\kappa_z|}{p_0}
=\beta_{\rm rel}^2\epsilon_{\rm win},
\end{align*}
其中用了 $m_\parallel=\gamma_0^3m_{\mathrm e}$ 与 $p_0=\gamma_0m_{\mathrm e} v_0$。因此 $\epsilon_{\rm win}$ 控制三阶相位相对于二阶相位的比例，而 $\varphi_3^{\max}$ 决定三阶绝对相位是否可分辨；二者承担不同角色。二阶反冲格在给定时间内要求 $\varphi_3^{\max}\ll1$。
\end{derivation}
